% Propietario conservado: /Users/ruben/Documents/New project/output/REVISION_ARGUMENTAL_APERTURAS_CIERRES_20260915/VI/delivery/MANUSCRIPT_VI_ES_EN_COMPLETE/source_es/sections/15_aritmetica_historias.tex
% Reunión focal desde I/sections/aritmetica_historias.tex, leído íntegro.
% Procedencia: technical/PROCEDENCIA_ARITMETICA_HISTORIAS_VI.md.
% Se conservan hist:seccion e hist:doble-varianza para el núcleo común.
% No incluye la representación de Weyl ni la prolongación excepcional de I.
\section{Arithmétique des histoires, composition et raffinement}
\label{vi:vi:hist:seccion-principal}

La persistance d’une particule et l’évaluation de son caractère nécessitent
de distinguer l’histoire transportée des coordonnées qui y sont lues.
Le noyau APP--TRIT--TPK fournit les états admissibles, leurs routes et
les opérations qui conservent leur provenance. Cette section précise trois
structures utilisées lors du passage de ces données à l’atlas : la compatibilité
entre profondeurs, la composition d’histoires avec mémoire et le
transport des observables en sens contraire au raffinement.
Les coefficients algébriques agissent sur des états et des routes déjà spécifiés ;
les relations prouvées ci-dessous conservent cet ordre.

\subsection{Sections compatibles et évaluation de leurs lecteurs}
\label{vi:hist:seccion}

Soit \(X_n\) l’ensemble des états admissibles de profondeur \(n\)
produits par le noyau. Un état enregistre les positions d’APP, les
feuillets additif et multiplicatif, leurs résidus et quotients, le régime
tritique, l’orientation, la phase, la frontière, la route, la retenue et la
mémoire pertinente pour la survie. Ces données satisfont les
relations des opérations qui les ont produites ; elles ne sont pas considérées comme un
produit de coordonnées indépendantes.

Pour \(m\ge n\), la restriction
\(r_{m,n}:X_m\longrightarrow X_n\) conserve l’antécédent de profondeur
\(n\). Les restrictions satisfont
\begin{equation}
 r_{n,n}=\operatorname{id}_{X_n},\qquad
 r_{\ell,n}=r_{m,n}\circ r_{\ell,m}
 \quad(\ell\ge m\ge n).
 \label{vi:vi:hist:restricciones}
\end{equation}
La nouvelle mémoire appartient au prolongement ; la restriction conserve celle
qui était déjà déterminée dans le préfixe.

\begin{definition}[Section numérique compatible]
Une section numérique HMT est une suite d’états de la même généalogie
\begin{equation}
 x=(x_n)_{n\ge0}\in X_\infty:=\varprojlim_n X_n,
 \qquad r_{m,n}(x_m)=x_n\quad(m\ge n).
 \label{vi:vi:hist:limite-inverso}
\end{equation}
Un lecteur est une opération définie sur un domaine déterminé de ces
sections. Une observation finie appartient à un préfixe ; une valeur scalaire
est l’évaluation d’un lecteur. La section et la présentation de cette section
par un lecteur sont donc des objets de types distincts.
\end{definition}

La notation de limite projective réunit la compatibilité des états déjà
construits. L’existence d’un prolongement concret est établie
par les règles d’admissibilité de cette construction, non par la
seule écriture de \(X_\infty\).

Dans le domaine archimédien d’un lecteur, chaque préfixe détermine un
intervalle fermé non vide \(I_n(x)=[a_n(x),b_n(x)]\), avec
\begin{equation}
 I_{n+1}(x)\subseteq I_n(x),\qquad
 b_n(x)-a_n(x)\longrightarrow0.
 \label{vi:vi:hist:cilindros}
\end{equation}
L’évaluation régionale par intervalles et la détermination de ses préfixes ont
été réunies dans la proposition~\ref{vi:vi:dep:prefijos}.

\begin{proposition}[Valeur déterminée par des cylindres compatibles]
Pour toute section de ce domaine, il existe un unique nombre réel
\(\operatorname{ev}(x)\) tel que
\begin{equation}
 \{\operatorname{ev}(x)\}=\bigcap_{n\ge0}I_n(x).
 \label{vi:vi:hist:evaluacion}
\end{equation}
\end{proposition}
\begin{proof}
L’emboîtement rend la suite \(a_n(x)\) croissante et
\(b_n(x)\) décroissante. Toutes deux restent bornées dans l’intervalle initial. Par
complétude de la droite, \(a=\sup_n a_n(x)\) et
\(b=\inf_n b_n(x)\) existent et satisfont \(a\le b\).
La largeur tendant vers zéro impose \(a=b\). Ce point appartient à chaque
intervalle et tout autre point de l’intersection aurait une distance à
celui-ci inférieure ou égale à toutes les largeurs ; il coïnciderait donc avec lui.
\end{proof}

Cette proposition détermine la valeur de chaque section du domaine indiqué.
Une couverture de tout \(\mathbb R\) exige de prouver que chaque valeur
appartient à l'image du lecteur ; la définition de limite projective
n'apporte pas à elle seule cette preuve. De même, l'identification de deux
histoires à partir de leur valeur exige un résultat d'injectivité
pour le lecteur considéré. L'évaluation conserve la généalogie comme
antécédent et spécifie exactement la partie qu'elle en publie.

\subsection{Algèbre sectorielle des routes et multiplicateurs de mémoire}
\label{vi:vi:hist:algebra}

Fixons un secteur d’états et la catégorie \(\mathcal C\) de leurs
histoires admissibles. Nous écrivons \(vu=v\circ u\) : on parcourt d’abord
\(u\), puis \(v\). Les pas élémentaires sont des opérations de
sélection, de transport et de mise à jour du TPK. La catégorie conserve
leurs extrémités et les conditions permettant de concaténer deux routes.

La fibre tritique locale peut être représentée par
\[
 \mathbb R[J_x]/(J_x^2+\tau(x)1_x),
 \qquad \tau(x)\in\{+1,0,-1\}.
\]
La représentation conserve le régime et son orientation. Pour traiter
des coefficients plus riches, nous admettrons une algèbre réelle associative et
unitale \(A_x\) en chaque objet \(x\). Une transition entre régimes
doit préciser son application sur ces algèbres ; elle ne revient pas à les identifier.

La construction suivante s’applique aux secteurs possédant des transports
unitaux \(\rho_u:A_x\to A_y\), pour \(u:x\to y\), avec
\begin{equation}
 \rho_{\operatorname{id}_x}=\operatorname{id}_{A_x},\qquad
 \rho_v\circ\rho_u=\rho_{vu}.
 \label{vi:vi:hist:accion}
\end{equation}
Pour chaque paire composable \(u:x\to y\), \(v:y\to z\), soit
\(\Omega(v,u)\in Z(A_z)^\times\) un multiplicateur central inversible.
Ses règles sont
\begin{equation}
 \Omega(\operatorname{id}_y,u)
 =\Omega(u,\operatorname{id}_x)=1_y
 \label{vi:vi:hist:normalizacion}
\end{equation}
et, pour \(w:z\to t\),
\begin{equation}
 \rho_w\!\left(\Omega(v,u)\right)\Omega(w,vu)
 =\Omega(w,v)\Omega(wv,u).
 \label{vi:vi:hist:cociclo}
\end{equation}
Ces conditions décrivent la structure de départ du secteur : chaque
application concrète doit fournir ses transports et ses multiplicateurs.
Une réalisation par bimodules conserve ses propres compositions et
n’est pas automatiquement remplacée par l’action stricte
\eqref{vi:vi:hist:accion}.

Soit \(t(u)\) l’extrémité terminale de la flèche \(u\). Sur les sommes
à support fini
\[
 \mathcal H=\bigoplus_{u\in\operatorname{Mor}(\mathcal C)}A_{t(u)}T_u
\]
nous définissons par bilinéarité
\begin{equation}
 (aT_v)(bT_u)=
 a\rho_v(b)\Omega(v,u)T_{vu}
 \label{vi:vi:hist:producto}
\end{equation}
lorsque \(u,v\) sont composables, et zéro sinon. La somme réunit
des histoires avec leurs coefficients ; le produit concatène les histoires et
conserve le multiplicateur de mémoire. Le support fini assure que chaque
produit est défini sans exiger de convergence supplémentaire.

\subsection{Associativité, identités locales et mémoire du calendrier}
\label{vi:vi:hist:asociatividad}

\begin{proposition}[Composition associative du secteur]
Sous \eqref{vi:vi:hist:accion}--\eqref{vi:vi:hist:cociclo}, le produit
\eqref{vi:vi:hist:producto} est associatif. S'il y a un nombre fini d'objets, son unité
est \(\sum_x1_xT_{\operatorname{id}_x}\). Pour un ensemble d'objets
arbitraire, les mêmes expressions sur des ensembles finis d'extrémités
donnent des identités locales pour toute famille finie d'histoires.
\end{proposition}
\begin{proof}
Soient \(u:x\to y\), \(v:y\to z\), \(w:z\to t\), avec des coefficients
\(a\in A_y\), \(b\in A_z\), \(c\in A_t\). Les deux associations
produisent comme coefficients de \(T_{wvu}\)
\begin{align*}
 ((cT_w)(bT_v))(aT_u):\quad&
 c\rho_w(b)\Omega(w,v)\rho_{wv}(a)\Omega(wv,u),\\
 (cT_w)((bT_v)(aT_u)):\quad&
 c\rho_w(b)\rho_w\rho_v(a)
 \rho_w(\Omega(v,u))\Omega(w,vu).
\end{align*}
L'action stricte identifie \(\rho_w\rho_v(a)\) à
\(\rho_{wv}(a)\). La centralité permet de déplacer
\(\Omega(w,v)\) à droite de ce facteur, et l'identité de cocycle
égalise les facteurs restants. Les coefficients généraux
\(a,b,c\), qui peuvent ne pas commuter, ne sont pas permutés. Si les trois flèches ne forment pas une
chaîne composable, les deux associations sont nulles par leurs extrémités.
La bilinéarité étend l'égalité aux sommes finies. Enfin,
\eqref{vi:vi:hist:normalizacion} et l'unitalité de chaque \(\rho_u\)
vérifient l'unité. Pour une famille finie, il suffit d'additionner les identités
de toutes ses extrémités initiales et terminales, ce qui démontre le dernier
énoncé.
\end{proof}

Le calendrier nonadique fournit une réalisation effective de ces
conditions. Prenons la catégorie à un objet avec des flèches
\(C_9=\mathbb Z/9\mathbb Z\), des coefficients
\(A=\mathbb R[z,z^{-1}]\) et l’action identité. Pour des représentants
\(0\le a,b<9\), nous définissons
\begin{equation}
 c_0(a,b)=\left\lfloor\frac{a+b}{9}\right\rfloor,
 \qquad\Omega(b,a)=z^{c_0(a,b)}.
 \label{vi:vi:hist:carry-calendario}
\end{equation}
La variable formelle \(z\) conserve le nombre de tours. La normalisation
découle de \(c_0(0,a)=c_0(a,0)=0\). Si \([a+b]_9\) désigne le
représentant réduit, la division euclidienne donne
\[
 c_0(a,b)+c_0([a+b]_9,c)
 =\left\lfloor\frac{a+b+c}{9}\right\rfloor
 =c_0(b,c)+c_0(a,[b+c]_9).
\]
Cette égalité prouve le cocycle \eqref{vi:vi:hist:cociclo} pour les
multiplicateurs, et donc l’associativité du calendrier avec mémoire.

Si \(T=T_{[1]_9}\), les règles du produit donnent
\(T^j=T_{[j]_9}\) pour \(0\le j<9\) et \(T^9=zT_{[0]_9}\).
En imposant ensuite \(z^{12}=1\), on obtient \(T^{108}=1\).
Les expressions \(z^kT_{[j]_9}\), \(0\le k<12\), \(0\le j<9\),
forment les \(108\) positions du calendrier résultant. On récupère ainsi
le groupe cyclique \(C_{108}\) et sa projection sur \(C_9\) ; conserver
\(z\) sans cette relation distingue les tours successifs. C’est une
réalisation de la mémoire de calendrier : les autres coordonnées des
histoires TPK restent dans leur domaine propre.

\subsection{Double variance du raffinement et conservation de l’unité}
\label{vi:hist:doble-varianza}

La restriction des états induit un prolongement des observables en sens
contraire. Pour les algèbres de fonctions munies du produit ponctuel
\(D_n=\operatorname{Fun}(X_n,\mathbb R)\), la précomposition définit
\begin{equation}
 r_{m,n}^*:D_n\longrightarrow D_m,
 \qquad r_{m,n}^*f=f\circ r_{m,n}.
 \label{vi:vi:hist:precomposicion}
\end{equation}
L’ordre contravariant
\(r_{\ell,n}^*=r_{\ell,m}^*\circ r_{m,n}^*\) est conservé. Ces applications
sont des homomorphismes unitaux. Si \(r_{m,n}\) est surjective, alors
\(r_{m,n}^*\) est également injective : une différence entre deux fonctions peut
être évaluée en un antécédent d’un point où elles diffèrent.

\begin{lemma}[Indicatrices, unité et évaluation compatible]
Pour \(x\in X_n\), soit \(e_x\) sa fonction indicatrice. Alors
\begin{equation}
 r_{m,n}^*e_x=\mathbf1_{r_{m,n}^{-1}(x)},\qquad
 r_{m,n}^*1_n=1_m.
 \label{vi:vi:hist:unidad}
\end{equation}
Lorsque la fibre \(r_{m,n}^{-1}(x)\) est finie, la première égalité est
la somme finie
\[
 r_{m,n}^*e_x=\sum_{y:\,r_{m,n}(y)=x}e_y.
\]
Pour toute section \(x=(x_n)\in X_\infty\), la règle
\([f]\mapsto f(x_n)\), avec \(f\in D_n\), définit un homomorphisme
unital
\begin{equation}
 \operatorname{Ev}_x:D_{\mathrm{cyl}}\longrightarrow\mathbb R,
 \qquad
 D_{\mathrm{cyl}}=\varinjlim_n(D_n,r_{n+1,n}^*).
 \label{vi:vi:hist:evaluacion-observables}
\end{equation}
\end{lemma}
\begin{proof}
Pour \(y\in X_m\), on a
\((r_{m,n}^*e_x)(y)=1\) exactement lorsque \(r_{m,n}(y)=x\).
C’est la première égalité de \eqref{vi:vi:hist:unidad}. Si la fibre est
finie, ses indicatrices ponctuelles sont orthogonales et leur somme est son indicatrice.
La seconde égalité s’obtient à partir de \(1_n(r_{m,n}(y))=1\), sans aucune
hypothèse sur le cardinal de la fibre.

Pour toute \(f\in D_n\), la naturalité de l’évaluation est
\[
 (r_{m,n}^*f)(y)=f(r_{m,n}(y)).
\]
Dans une section compatible, \(r_{m,n}(x_m)=x_n\), donc
\((r_{m,n}^*f)(x_m)=f(x_n)\). Par conséquent, deux représentants qui
s’identifient à une profondeur commune ont la même évaluation.
L’évaluation ponctuelle préserve somme, produit et unité, ce qui conclut
la preuve de \eqref{vi:vi:hist:evaluacion-observables}.
\end{proof}

Pour les fibres infinies, la formulation au moyen de
\(\mathbf1_{r_{m,n}^{-1}(x)}\) reste valide ; il n’est pas nécessaire d’interpréter une somme
infinie comme une somme algébrique finie. La limite inductive réunit les
observables de profondeur finie avec les identifications induites par
les restrictions. Le prolongement des opérateurs de route exige, en outre,
les carrés de compatibilité de leurs transports avec ces restrictions.

L’unité conservée est la fonction constante égale à un sur le domaine des
possibilités. Une section individuelle sélectionne une histoire compatible
de ce domaine. Le raffinement peut augmenter la profondeur et distinguer
davantage de prolongements tout en préservant l’unité et l’évaluation de chaque
antécédent. Cette conservation unitale fournit le contenu algébrique
précis utilisé ici ; une mesure de probabilité, lorsqu’elle intervient dans
un autre lecteur, nécessite en outre sa normalisation et ses propres règles.
